\documentclass[10pt,a4paper]{amsart}
\usepackage{amsmath,amssymb,amsthm,graphicx}
\usepackage{verbatim}
\newenvironment{reference}{\comment}{\endcomment}
%\newenvironment{reference}{}{}
\newenvironment{slogan}{\comment}{\endcomment}
\newenvironment{history}{\comment}{\endcomment}

% For commutative diagrams we use Xy-pic
\usepackage[all]{xy}

% We use 2cell for 2-commutative diagrams.
\xyoption{2cell}
% Theorem environments.
%
\theoremstyle{plain}
\newtheorem{theorem}[subsection]{Theorem}
\newtheorem{proposition}[subsection]{Proposition}
\newtheorem{lemma}[subsection]{Lemma}

\theoremstyle{definition}
\newtheorem{definition}[subsection]{Definition}
\newtheorem{example}[subsection]{Example}
\newtheorem{exercise}[subsection]{Exercise}
\newtheorem{situation}[subsection]{Situation}

\theoremstyle{remark}
\newtheorem{remark}[subsection]{Remark}
\newtheorem{remarks}[subsection]{Remarks}

\numberwithin{equation}{subsection}

% Macros
%
\def\lim{\mathop{\mathrm{lim}}\nolimits}
\def\colim{\mathop{\mathrm{colim}}\nolimits}
\def\Spec{\mathop{\mathrm{Spec}}}
\def\Hom{\mathop{\mathrm{Hom}}\nolimits}
\def\Ext{\mathop{\mathrm{Ext}}\nolimits}
\def\SheafHom{\mathop{\mathcal{H}\!\mathit{om}}\nolimits}
\def\SheafExt{\mathop{\mathcal{E}\!\mathit{xt}}\nolimits}
\def\Sch{\mathit{Sch}}
\def\Mor{\mathop{\mathrm{Mor}}\nolimits}
\def\Ob{\mathop{\mathrm{Ob}}\nolimits}
\def\Sh{\mathop{\mathit{Sh}}\nolimits}
\def\NL{\mathop{N\!L}\nolimits}
\def\CH{\mathop{\mathrm{CH}}\nolimits}
\def\proetale{{pro\text{-}\acute{e}tale}}
\def\etale{{\acute{e}tale}}
\def\QCoh{\mathit{QCoh}}
\def\Ker{\mathop{\mathrm{Ker}}}
\def\Im{\mathop{\mathrm{Im}}}
\def\Coker{\mathop{\mathrm{Coker}}}
\def\Coim{\mathop{\mathrm{Coim}}}

% Boxtimes
%
\DeclareMathSymbol{\boxtimes}{\mathbin}{AMSa}{"02}

%
% Macros for moduli stacks/spaces
%
\def\QCohstack{\mathcal{QC}\!\mathit{oh}}
\def\Cohstack{\mathcal{C}\!\mathit{oh}}
\def\Spacesstack{\mathcal{S}\!\mathit{paces}}
\def\Quotfunctor{\mathrm{Quot}}
\def\Hilbfunctor{\mathrm{Hilb}}
\def\Curvesstack{\mathcal{C}\!\mathit{urves}}
\def\Polarizedstack{\mathcal{P}\!\mathit{olarized}}
\def\Complexesstack{\mathcal{C}\!\mathit{omplexes}}
% \Pic is the operator that assigns to X its picard group, usage \Pic(X)
% \Picardstack_{X/B} denotes the Picard stack of X over B
% \Picardfunctor_{X/B} denotes the Picard functor of X over B
\def\Pic{\mathop{\mathrm{Pic}}\nolimits}
\def\Picardstack{\mathcal{P}\!\mathit{ic}}
\def\Picardfunctor{\mathrm{Pic}}
\def\Deformationcategory{\mathcal{D}\!\mathit{ef}}
\usepackage{cleveref}
\crefname{originalreference}{original reference}{original references}
\Crefname{originalreference}{Original reference}{Original references}
\makeatletter
\newlabel{lemma-extend-over-dvr}{{0BXY}{source}{}{sourcepointer-1}{}}
\newlabel{lemma-extend-over-dvr@cref}{{[originalreference][0][]0BXY}{[1][0][]source}{}{}{}}
\makeatother
\crefname{originalreference}{original reference}{original references}
\Crefname{originalreference}{Original reference}{Original references}
\makeatletter
\newlabel{lemma-extend-over-normal-curve}{{0BXZ}{source}{}{sourcepointer-2}{}}
\newlabel{lemma-extend-over-normal-curve@cref}{{[originalreference][0][]0BXZ}{[1][0][]source}{}{}{}}
\makeatother
\crefname{originallemma}{original lemma}{original lemmas}
\Crefname{originallemma}{Original lemma}{Original lemmas}
\makeatletter
\newlabel{morphisms-lemma-relative-normalization-normal-codim-1}{{0BXB}{source}{}{sourcepointer-3}{}}
\newlabel{morphisms-lemma-relative-normalization-normal-codim-1@cref}{{[originallemma][0][]0BXB}{[1][0][]source}{}{}{}}
\makeatother
\crefname{originaldefinition}{original definition}{original definitions}
\Crefname{originaldefinition}{Original definition}{Original definitions}
\makeatletter
\newlabel{properties-definition-regular}{{02IS}{source}{}{sourcepointer-4}{}}
\newlabel{properties-definition-regular@cref}{{[originaldefinition][0][]02IS}{[1][0][]source}{}{}{}}
\makeatother
\crefname{originallemma}{original lemma}{original lemmas}
\Crefname{originallemma}{Original lemma}{Original lemmas}
\makeatletter
\newlabel{properties-lemma-normal-dimension-1-regular}{{0BX2}{source}{}{sourcepointer-5}{}}
\newlabel{properties-lemma-normal-dimension-1-regular@cref}{{[originallemma][0][]0BX2}{[1][0][]source}{}{}{}}
\makeatother
\crefname{originallemma}{original lemma}{original lemmas}
\Crefname{originallemma}{Original lemma}{Original lemmas}
\makeatletter
\newlabel{properties-lemma-regular-normal}{{0569}{source}{}{sourcepointer-6}{}}
\newlabel{properties-lemma-regular-normal@cref}{{[originallemma][0][]0569}{[1][0][]source}{}{}{}}
\makeatother
\crefname{originallemma}{original lemma}{original lemmas}
\Crefname{originallemma}{Original lemma}{Original lemmas}
\makeatletter
\newlabel{varieties-lemma-dim-1-projective-completion}{{0BXV}{source}{}{sourcepointer-7}{}}
\newlabel{varieties-lemma-dim-1-projective-completion@cref}{{[originallemma][0][]0BXV}{[1][0][]source}{}{}{}}
\makeatother
\crefname{originallemma}{original lemma}{original lemmas}
\Crefname{originallemma}{Original lemma}{Original lemmas}
\makeatletter
\newlabel{varieties-lemma-dim-1-proper-projective}{{0A26}{source}{}{sourcepointer-8}{}}
\newlabel{varieties-lemma-dim-1-proper-projective@cref}{{[originallemma][0][]0A26}{[1][0][]source}{}{}{}}
\makeatother
\crefname{originallemma}{original lemma}{original lemmas}
\Crefname{originallemma}{Original lemma}{Original lemmas}
\makeatletter
\newlabel{varieties-lemma-dimension-locally-algebraic}{{0A21}{source}{}{sourcepointer-9}{}}
\newlabel{varieties-lemma-dimension-locally-algebraic@cref}{{[originallemma][0][]0A21}{[1][0][]source}{}{}{}}
\makeatother
\crefname{originallemma}{original lemma}{original lemmas}
\Crefname{originallemma}{Original lemma}{Original lemmas}
\makeatletter
\newlabel{varieties-lemma-normalization-locally-algebraic}{{0BXR}{source}{}{sourcepointer-10}{}}
\newlabel{varieties-lemma-normalization-locally-algebraic@cref}{{[originallemma][0][]0BXR}{[1][0][]source}{}{}{}}
\makeatother
\crefname{originallemma}{original lemma}{original lemmas}
\Crefname{originallemma}{Original lemma}{Original lemmas}
\makeatletter
\newlabel{varieties-lemma-reduced-dim-1-projective-completion}{{0BXW}{source}{}{sourcepointer-11}{}}
\newlabel{varieties-lemma-reduced-dim-1-projective-completion@cref}{{[originallemma][0][]0BXW}{[1][0][]source}{}{}{}}
\makeatother
\crefname{originallemma}{original lemma}{original lemmas}
\Crefname{originallemma}{Original lemma}{Original lemmas}
\makeatletter
\newlabel{varieties-lemma-regular-point-on-curve}{{0B8Y}{source}{}{sourcepointer-12}{}}
\newlabel{varieties-lemma-regular-point-on-curve@cref}{{[originallemma][0][]0B8Y}{[1][0][]source}{}{}{}}
\makeatother
\crefname{originaltheorem}{original theorem}{original theorems}
\Crefname{originaltheorem}{Original theorem}{Original theorems}
\makeatletter
\newlabel{varieties-theorem-varieties-rational-maps}{{0BXN}{source}{}{sourcepointer-13}{}}
\newlabel{varieties-theorem-varieties-rational-maps@cref}{{[originaltheorem][0][]0BXN}{[1][0][]source}{}{}{}}
\makeatother
\title{Canonical equivalence of function fields of transcendence degree 1 with curves under dominant rational maps and with normal projective curves under morphisms}
\author{The Stacks Project Authors}
\date{stacks-project:curves.tex:local-snapshot-20261012; GNU FDL 1.2 or later}
\begin{document}
\maketitle
\noindent\textit{Editorial note.} All delivered mathematical text is the author's own, copied from the pinned Stacks Project chapter file curves.tex; no statement or proof step was rewritten, reordered or removed.\par\medskip
\noindent\textit{Editorial note.} The counted proof is the author's complete argument: the restriction of the varieties-level equivalence, the regular-normal-projective dictionary, the fully faithful functor with its two extension lemmas, the normal projective model construction, and the author's closing remark on the alternative compactify-then-normalize order are all retained verbatim.\par\medskip
\noindent\textit{Editorial note.} The label theorem-curves-rational-maps names the counted statement, and the proof environment immediately following it in the source is delivered verbatim as its proof.\par\medskip
\noindent\textit{Editorial note.} The counted statement is carried by the author's own theorem environment in the chapter source.\par\medskip
\noindent\textit{Allowed standard prerequisites (editorial).}\begin{itemize}
\item Rational maps of varieties and function fields. Varieties, Theorem varieties-theorem-varieties-rational-maps with Lemma varieties-lemma-dimension-locally-algebraic. Supplies the equivalence between (1) and (2) by restriction: a variety is a curve exactly when its function field has transcendence degree 1.
\item Extension of morphisms from the generic point. Same chapter, Lemmas lemma-extend-over-normal-curve and lemma-extend-over-dvr; Varieties, Lemmas varieties-lemma-dim-1-proper-projective, varieties-lemma-normalization-locally-algebraic, varieties-lemma-reduced-dim-1-projective-completion. Extends a rational map from a normal curve to a morphism, and produces the normal projective model of a given function field by normalizing and completing.
\end{itemize}\medskip
\ifdefined\hypertarget\hypertarget{sourcepointer-1}{}\fi
\noindent\textit{Source reference.} Original reference 0BXY: sources/stacks-curves/pinned-original-source.tex, source lines 272--272. Prerequisite pointer invoked at proof line(s) 421; it names a standard fact of the same chapter that the counted argument uses, and every step of the counted proof built on it is retained verbatim in the delivered body.\par\medskip
\ifdefined\hypertarget\hypertarget{sourcepointer-2}{}\fi
\noindent\textit{Source reference.} Original reference 0BXZ: sources/stacks-curves/pinned-original-source.tex, source lines 287--287. Prerequisite pointer invoked at proof line(s) 418; it names a standard fact of the same chapter that the counted argument uses, and every step of the counted proof built on it is retained verbatim in the delivered body.\par\medskip
\ifdefined\hypertarget\hypertarget{sourcepointer-3}{}\fi
\noindent\textit{Source reference.} Original lemma 0BXB: sources/stacks-chow/tags.txt, source lines 14974--14974. Prerequisite pointer invoked at proof line(s) 444; it names a standard fact of a sibling chapter that the counted argument uses, and every step of the counted proof built on it is retained verbatim in the delivered body.\par\medskip
\ifdefined\hypertarget\hypertarget{sourcepointer-4}{}\fi
\noindent\textit{Source reference.} Original definition 02IS: sources/stacks-chow/tags.txt, source lines 3311--3311. Prerequisite pointer invoked at proof line(s) 400; it names a standard fact of a sibling chapter that the counted argument uses, and every step of the counted proof built on it is retained verbatim in the delivered body.\par\medskip
\ifdefined\hypertarget\hypertarget{sourcepointer-5}{}\fi
\noindent\textit{Source reference.} Original lemma 0BX2: sources/stacks-chow/tags.txt, source lines 14965--14965. Prerequisite pointer invoked at proof line(s) 404; it names a standard fact of a sibling chapter that the counted argument uses, and every step of the counted proof built on it is retained verbatim in the delivered body.\par\medskip
\ifdefined\hypertarget\hypertarget{sourcepointer-6}{}\fi
\noindent\textit{Source reference.} Original lemma 0569: sources/stacks-chow/tags.txt, source lines 6702--6702. Prerequisite pointer invoked at proof line(s) 403; it names a standard fact of a sibling chapter that the counted argument uses, and every step of the counted proof built on it is retained verbatim in the delivered body.\par\medskip
\ifdefined\hypertarget\hypertarget{sourcepointer-7}{}\fi
\noindent\textit{Source reference.} Original lemma 0BXV: sources/stacks-chow/tags.txt, source lines 14993--14993. Prerequisite pointer invoked at proof line(s) 440; it names a standard fact of a sibling chapter that the counted argument uses, and every step of the counted proof built on it is retained verbatim in the delivered body.\par\medskip
\ifdefined\hypertarget\hypertarget{sourcepointer-8}{}\fi
\noindent\textit{Source reference.} Original lemma 0A26: sources/stacks-chow/tags.txt, source lines 12694--12694. Prerequisite pointer invoked at proof line(s) 408; it names a standard fact of a sibling chapter that the counted argument uses, and every step of the counted proof built on it is retained verbatim in the delivered body.\par\medskip
\ifdefined\hypertarget\hypertarget{sourcepointer-9}{}\fi
\noindent\textit{Source reference.} Original lemma 0A21: sources/stacks-chow/tags.txt, source lines 12689--12689. Prerequisite pointer invoked at proof line(s) 395; it names a standard fact of a sibling chapter that the counted argument uses, and every step of the counted proof built on it is retained verbatim in the delivered body.\par\medskip
\ifdefined\hypertarget\hypertarget{sourcepointer-10}{}\fi
\noindent\textit{Source reference.} Original lemma 0BXR: sources/stacks-chow/tags.txt, source lines 14989--14989. Prerequisite pointer invoked at proof line(s) 432; it names a standard fact of a sibling chapter that the counted argument uses, and every step of the counted proof built on it is retained verbatim in the delivered body.\par\medskip
\ifdefined\hypertarget\hypertarget{sourcepointer-11}{}\fi
\noindent\textit{Source reference.} Original lemma 0BXW: sources/stacks-chow/tags.txt, source lines 14994--14994. Prerequisite pointer invoked at proof line(s) 435; it names a standard fact of a sibling chapter that the counted argument uses, and every step of the counted proof built on it is retained verbatim in the delivered body.\par\medskip
\ifdefined\hypertarget\hypertarget{sourcepointer-12}{}\fi
\noindent\textit{Source reference.} Original lemma 0B8Y: sources/stacks-chow/tags.txt, source lines 14156--14156. Prerequisite pointer invoked at proof line(s) 405; it names a standard fact of a sibling chapter that the counted argument uses, and every step of the counted proof built on it is retained verbatim in the delivered body.\par\medskip
\ifdefined\hypertarget\hypertarget{sourcepointer-13}{}\fi
\noindent\textit{Source reference.} Original theorem 0BXN: sources/stacks-chow/tags.txt, source lines 14986--14986. Prerequisite pointer invoked at proof line(s) 392; it names a standard fact of a sibling chapter that the counted argument uses, and every step of the counted proof built on it is retained verbatim in the delivered body.\par\medskip
\begin{theorem}
\label{theorem-curves-rational-maps}
Let $k$ be a field. The following categories are canonically equivalent
\begin{enumerate}
\item The category of finitely generated field extensions $K/k$ of
transcendence degree $1$.
\item The category of curves and dominant rational maps.
\item The category of normal projective curves and nonconstant morphisms.
\item The category of nonsingular projective curves and nonconstant morphisms.
\item The category of regular projective curves and nonconstant morphisms.
\item The category of normal proper curves and nonconstant morphisms.
\end{enumerate}
\end{theorem}

\begin{proof}
The equivalence between categories (1) and (2) is the restriction of the
equivalence of
Varieties, Theorem \ref{varieties-theorem-varieties-rational-maps}.
Namely, a variety is a curve if and only if its function field has
transcendence degree $1$, see for example
Varieties, Lemma \ref{varieties-lemma-dimension-locally-algebraic}.

\medskip\noindent
The categories in (3), (4), (5), and (6) are the same. First of all, the
terms ``regular'' and ``nonsingular'' are synonyms, see
Properties, Definition \ref{properties-definition-regular}.
Being normal and regular are the same thing for Noetherian
$1$-dimensional schemes
(Properties, Lemmas \ref{properties-lemma-regular-normal} and
\ref{properties-lemma-normal-dimension-1-regular}). See
Varieties, Lemma \ref{varieties-lemma-regular-point-on-curve}
for the case of curves. Thus (3) is the same as (5). Finally, (6)
is the same as (3) by
Varieties, Lemma \ref{varieties-lemma-dim-1-proper-projective}.

\medskip\noindent
If $f : X \to Y$ is a nonconstant morphism of nonsingular projective curves,
then $f$ sends the generic point $\eta$ of $X$ to the generic point $\xi$ of
$Y$. Hence we obtain a morphism
$k(Y) = \mathcal{O}_{Y, \xi} \to \mathcal{O}_{X, \eta} = k(X)$
in the category (1). If two morphisms $f,g: X \to Y$ gives the same morphism
$k(Y) \to k(X)$, then by the equivalence between (1) and (2),
$f$ and $g$ are equivalent as rational maps, so $f=g$ by
Lemma \ref{lemma-extend-over-normal-curve}.
Conversely, suppose that we have a map
$k(Y) \to k(X)$ in the category (1). Then we obtain a morphism $U \to Y$
for some nonempty open $U \subset X$. By Lemma \ref{lemma-extend-over-dvr}
this extends to all of $X$ and we obtain a morphism in the category (5).
Thus we see that there is a fully faithful functor (5)$\to$(1).

\medskip\noindent
To finish the proof we have to show that every $K/k$ in (1)
is the function field of a normal projective curve.
We already know that $K = k(X)$ for some curve $X$.
After replacing $X$ by its normalization
(which is a variety birational to $X$)
we may assume $X$ is normal
(Varieties, Lemma \ref{varieties-lemma-normalization-locally-algebraic}).
Then we choose $X \to \overline{X}$ with
$\overline{X} \setminus X = \{x_1, \ldots, x_n\}$ as in
Varieties, Lemma \ref{varieties-lemma-reduced-dim-1-projective-completion}.
Since $X$ is normal and since each
of the local rings $\mathcal{O}_{\overline{X}, x_i}$ is normal
we conclude that $\overline{X}$ is a normal projective curve as desired.
(Remark: We can also first compactify using
Varieties, Lemma \ref{varieties-lemma-dim-1-projective-completion}
and then normalize using
Varieties, Lemma \ref{varieties-lemma-normalization-locally-algebraic}.
Doing it this way we avoid using the somewhat tricky
Morphisms, Lemma \ref{morphisms-lemma-relative-normalization-normal-codim-1}.)
\end{proof}

\end{document}
